%%% Local Variables:
%%% mode: latex
%%% TeX-master: "../doc"
%%% coding: utf-8
%%% End:
% !TEX TS-program = pdflatexmk
% !TEX encoding = UTF-8 Unicode
% !TEX root = ../doc.tex

This chapter presents the findings of the user study conducted to evaluate how different
explainability methods affect user trust in large language model (LLM) outputs. The study
compared five conditions: a plain LLM baseline, Chain-of-Thought reasoning (CoT), SHAP
feature attributions, confidence scores (CONF), and an explanation on the working of an LLM.
Trust was measured across five Likert-scaled items (1\,=\,strongly disagree,
5\,=\,strongly agree). The following sections first characterise the participant sample,
then report trust outcomes per method, and finally describe the relative helpfulness
rankings provided by participants.

\section{Sample Characteristics}
\label{sec:sample}
 
\subsection{Familiarity and Prior Use of Generative AI}
 
Of all survey respondents, \textbf{67.3\%} reported being familiar with generative AI
tools such as ChatGPT, Gemini, Copilot, Claude, LLaMA, or Bard (see
Figure~\ref{fig:familiar}). of those familiar with genAI \textbf{77.4\%} indicated having
used at least one such tool at least once (Figure~\ref{fig:used}). Among those who had
never used an LLM, the most frequently cited barrier was \emph{not yet having seen a
need} ($n=37$), followed by a preference for traditional search engines ($n=21$) and
privacy or security concerns ($n=17$) (Figure~\ref{fig:nonusers}).
 
\begin{figure}[ht]
  \centering
  % \includegraphics[width=0.5\textwidth]{figures/familiar_pie.png}
  \caption{Self-reported familiarity with generative AI tools.}
  \label{fig:familiar}
\end{figure}
 
\begin{figure}[ht]
  \centering
  % \includegraphics[width=0.5\textwidth]{figures/used_pie.png}
  \caption{Self-reported prior use of a generative AI tool at least once.}
  \label{fig:used}
\end{figure}
 
\begin{figure}[ht]
  \centering
  % \includegraphics[width=0.9\textwidth]{figures/nonusers_bar.png}
  \caption{Reasons given by participants who had never used a large language model.}
  \label{fig:nonusers}
\end{figure}
 
Among LLM users, ChatGPT (OpenAI) was by far the most widely tried model ($n \approx 150$),
followed by Gemini (Google DeepMind, $n \approx 85$) and Microsoft Copilot ($n \approx 51$)
(Figure~\ref{fig:llms_tried}). Claude (Anthropic) had been used by approximately 12
participants. This distribution indicates that the sample had broad, realistic exposure to
current commercial LLMs, though it was skewed toward the most mainstream tools.
 
\begin{figure}[ht]
  \centering
  % \includegraphics[width=\textwidth]{figures/llms_tried_bar.png}
  \caption{LLMs tried by participants, sorted by usage count.}
  \label{fig:llms_tried}
\end{figure}
 
\subsection{Employment Role and Sector}
\label{sec:sample_demographics}
 
Figure~\ref{fig:role} shows LLM usage broken down by employment role. Administration
($n=26$), Retired ($n=21$), and Specialist/Expert ($n=19$) were the three most represented
groups, while technical roles such as IT/Software/Data ($n=15$) and Technical
Staff/Engineer ($n=11$) together comprised a moderate share of users. The broad spread
across non-technical roles suggests that the findings are not exclusively representative
of a technically literate audience.
 
\begin{figure}[ht]
  \centering
  % \includegraphics[width=\textwidth]{figures/role_bar.png}
  \caption{Number of LLM users per employment role.}
  \label{fig:role}
\end{figure}
 
Sector-level results (Figure~\ref{fig:sector}) show Healthcare/Pharmaceuticals/Biotechnology
as the largest group ($n \approx 26$), followed by participants who did not specify a sector
($n \approx 22$) and Public Sector/Administration ($n \approx 18$). Education/Research
($n \approx 17$) also contributed substantially.
 
These demographic patterns introduce a potential \textbf{response bias}: the over-representation of administration, healthcare, and retirement demographics---relative to, say, software engineering---may lead to conservative trust ratings and a stronger preference for intuitive, human-readable explanations (such as CoT) over technical ones (such as SHAP). This limitation is discussed further in Chapter~\ref{chap:discussion}.
 
\begin{figure}[ht]
  \centering
  % \includegraphics[width=\textwidth]{figures/sector_bar.png}
  \caption{Number of LLM users per industry sector.}
  \label{fig:sector}
\end{figure}
 
\subsection{Education Level}
 
LLM usage was highest among participants with a \emph{Realschulabschluss} (intermediate
secondary certificate, $n=54$) and among those holding an \emph{Abitur}
(university-entrance qualification, $n=36$), followed by Bachelor's ($n=29$) and
Master's degree holders ($n=28$) (Figure~\ref{fig:edu}). The relatively high LLM adoption
at the secondary rather than university level is notable and may partly reflect the larger
absolute population of secondary-level respondents in the panel.
 
\begin{figure}[ht]
  \centering
  % \includegraphics[width=0.9\textwidth]{figures/edu_bar.png}
  \caption{LLM usage per education level (German qualification labels).}
  \label{fig:edu}
\end{figure}

Diagrams and stuff showing results from survey
- COT is the most impactful (quite closely followed by confidence score)
    - the most people ranked it as the most useful and changed their opinion on all trust aspects the most
    - also the lowest overall in trust
- research other methods in this category human centered explanations




Discuss bias: 
- llm usage by sector and job => what people the pannel consists of




and explaining the resulting prototype implementation

\section{Prototype Implementation}
The resulting best xai method has been implemented like this

- implement a prototype of COT combined with CONF

